\par\medskip
\noindent
13.4. \underline{Interpretación mediante las} «\underline{estructuras de Hodge
mixtas}» \underline{de Deligne.}

Volvamos primero al diccionario de 13.1.4, que reformularemos de una
manera ligeramente distinta, más adecuada para introducir las
«estructuras de Hodge» pertinentes en la situación presente. Para ello,
sea $\underline R$ el haz de secciones de $R$ sobre $S$, e introduzcamos
el $\underline{\mathcal O}_S$-módulo

\begin{equation}\tag{13.4.1}
 \mathcal E(A)=\mathcal E=
 \underline R\otimes_{\mathbb Z}\underline{\mathcal O}_S\qquad .
\end{equation}

\noindent
Es localmente libre de tipo finito cuando $R$ es localmente constante.
(Obsérvese que $\mathcal E$ no es necesariamente coherente cuando $R$
no es localmente constante.) El homomorfismo (13.1.3)
$\underline R\longrightarrow\underline E$ define entonces un
homomorfismo canónico

\begin{equation}\tag{13.4.2}
 \mathcal E=\underline R\otimes_{\mathbb Z}
 \underline{\mathcal O}_S\longrightarrow\underline E\qquad .
\end{equation}

\noindent
Denotamos su núcleo por $\underline F(A)$, o simplemente por
$\underline F$; se tiene así una sucesión exacta

\begin{equation}\tag{13.4.3}
 0\longrightarrow\underline F\longrightarrow\mathcal E
 \longrightarrow\underline E\qquad .
\end{equation}

\noindent
El caso más interesante es aquel en que (13.4.2) es un epimorfismo, es
decir, aquel en que $R_s$ genera el espacio vectorial complejo $E_s$
para todo $s\in S$. Se tiene entonces una sucesión exacta

\begin{equation}\tag{13.4.4}
 0\longrightarrow\underline F\longrightarrow\mathcal E
 \longrightarrow\underline E\longrightarrow0\qquad .
\end{equation}

\noindent
Si $R$ es localmente constante, esta es una sucesión exacta corta de
módulos localmente libres sobre $S$. Conocer el par $(E,R)$ equivale
entonces a conocer el par $(R,\underline F)$, donde $R$ es un grupo
analítico localmente constante sobre $S$ cuyas fibras son
$\mathbb Z$-módulos libres de tipo finito, y $\underline F$ es un
submódulo localmente sumando directo de
$\mathcal E=\underline R\otimes_{\mathbb Z}
\underline{\mathcal O}_S$, sometido además a

\begin{equation}\tag{13.4.5}
 \underline F\cap\underline R=0\qquad .
\end{equation}

\noindent
Desde el punto de vista de las estructuras de Hodge, $\mathcal E$ debe
considerarse filtrado por la «\underline{filtración de Hodge}»

\begin{equation}\tag{13.4.6}
 \operatorname{Fil}^{\sisi{1}{-1}}(\mathcal E)=\mathcal E,
 \qquad \operatorname{Fil}^{0}(\mathcal E)=\underline F,
 \qquad \operatorname{Fil}^{\sisi{-1}{1}}(\mathcal E)=0\qquad .
\end{equation}

\noindent
Obsérvese que la sucesión exacta (13.4.3) sigue dependiendo
funtorialmente del grupo analítico $A$ sobre $S$, y es compatible con
cualquier cambio de base.

Volviendo ahora a las condiciones de 13.3, pongamos

\begin{equation}\tag{13.4.7}
 \mathcal E^f=\underline R^f\otimes_{\mathbb Z}
 \underline{\mathcal O}_S,
 \qquad
 \mathcal E^t=\underline R^t\otimes_{\mathbb Z}
 \underline{\mathcal O}_S\qquad .
\end{equation}

\noindent
Estos módulos están definidos sobre un entorno adecuado de $s$, que
podemos suponer igual a todo $S$, y definen una filtración de
$\mathcal E$, llamada \underline{filtración por pesos},

\begin{equation}\tag{13.4.8}
 \mathcal E\supset\mathcal E^f\supset\mathcal E^t\supset0\qquad .
\end{equation}

\noindent
Aquí $\mathcal E^t$, $\mathcal E^f$ y $\mathcal E$ deben considerarse
de pesos $-2$, $-1$ y $0$, respectivamente; se trata, por tanto, de una
filtración creciente. Al igual que $R^t$, el módulo $\mathcal E^t$
depende de la elección de un sub-toro $T_o$ de $A_o=A_s$. Puesto que
$R^f$ y $R^t$ son localmente constantes, \underline{$\mathcal E^f$ y
$\mathcal E^t$ son localmente libres}; también pueden interpretarse
como

\begin{equation}\tag{13.4.9}
 \mathcal E^f=\mathcal E(A^\natural),
 \qquad \mathcal E^t=\mathcal E(T)\qquad .
\end{equation}

\noindent
Aplicando la functorialidad de (13.4.3), cuando varía $A$, a los
homomorfismos

\begin{equation*}
 T\longrightarrow A^\natural\longrightarrow A
\end{equation*}

\noindent
se ve inmediatamente que $\underline F(A^\natural)$ y $\underline F(T)$
se identifican, respectivamente, con

\begin{equation}\tag{13.4.10}
 \underline F(A^\natural)\simeq
 \underline F^f\stackrel{\mathrm{def}}{=}
 \underline F\cap\mathcal E^f,
 \qquad
 \underline F(T)=\underline F^t=
 \underline F\cap\mathcal E^t=0\qquad .
\end{equation}

\noindent
La última igualdad se deduce inmediatamente de que
$\underline F(T)=0$ para un toro $T$: en este caso, la sucesión exacta
de tipo (13.4.3) es
$0\longrightarrow0\longrightarrow\mathcal E(T)
\longrightarrow\underline{\operatorname{Lie}}(T)\longrightarrow0$.
Al mismo tiempo se ve que el homomorfismo compuesto

\begin{equation*}
 \mathcal E^t\longrightarrow\mathcal E\longrightarrow\underline E
\end{equation*}

\noindent
induce un isomorfismo

\begin{equation}\tag{13.4.11}
 \mathcal E^t\stackrel{\sim}{\longrightarrow}\underline E^t
 \simeq\underline{\operatorname{Lie}}(T)\qquad .
\end{equation}

\noindent
\underline{Supongamos ahora que $R_s$ genera el espacio vectorial
complejo $E_s$}, como sucede en particular en el caso favorable 13.3.16,
en el que $A_s$ es una extensión de un grupo abeloide por un toro
complejo $T_o$. Tras restringir $S$ a un entorno adecuado de $s$, el
homomorfismo canónico de tipo (13.4.2), con $A$ sustituido por
$A^\natural$,

\begin{equation*}
 \mathcal E^f=\mathcal E(A^\natural)
 \longrightarrow\underline{\operatorname{Lie}}(A^\natural)
 \simeq\underline E
\end{equation*}

\noindent
es un epimorfismo; con mayor razón, también lo es (13.4.2). Tenemos,
pues, las sucesiones exactas (13.4.4) y

\begin{equation}\tag{13.4.12}
 0\longrightarrow\underline F^f\longrightarrow\mathcal E^f
 \longrightarrow\underline E\longrightarrow0\qquad .
\end{equation}

\noindent
En particular,

\begin{equation}\tag{13.4.13}
 \underline F+\mathcal E^f=\mathcal E\qquad ,
\end{equation}

\noindent
pues $\mathcal E/\underline F\simeq\underline E$. Además, como
$B=A^\natural/T$ es un cociente de \sisi{$A$}{$A^\natural$},
\footnote{La fuente impresa afirma aquí que $B$ es un cociente de $A$;
la definición mostrada $B=A^\natural/T$ determina el grupo pretendido.}
satisface la condición de sobreyectividad análoga y da, por tanto,
lugar a una sucesión exacta de tipo (13.4.4):

\begin{equation}\tag{13.4.14}
 0\longrightarrow\underline F(B)\longrightarrow\mathcal E(B)
 \longrightarrow\underline{\operatorname{Lie}}(B)
 \longrightarrow0\qquad .
\end{equation}

\noindent
Su interpretación mediante los módulos ya introducidos es inmediata.
En efecto, por (13.3.12),

\begin{equation}\tag{13.4.15}
 \left\{
 \begin{array}{l}
  \mathcal E(B)\simeq
  (\underline R^f/\underline R^t)\otimes_{\mathbb Z}
  \underline{\mathcal O}_S
  \simeq\mathcal E^f/\mathcal E^t,\\[2ex]
  \underline{\operatorname{Lie}}(B)
  \simeq\underline E/\underline E^t.
 \end{array}
 \right.
\end{equation}

\noindent
De (13.4.11) y de la relación
$\underline F\cap\mathcal E^t=0$ de (13.4.10) se sigue que
$\underline F(B)=\ker(\mathcal E^f/\mathcal E^t
\longrightarrow\underline E/\underline E^t)$ es isomorfo a
$\underline F^f$ mediante el homomorfismo

\begin{equation}\tag{13.4.16}
 \underline F^f\simeq\underline F(A^\natural)
 \stackrel{\sim}{\longrightarrow}\underline F(B)\qquad .
\end{equation}

\noindent
Resumamos las propiedades que acabamos de obtener para la filtración
por pesos (13.4.8), y su relación con la filtración de Hodge (13.4.6)
de $\mathcal E$ definida por el submódulo $\underline F$, en el caso
favorable 13.3.16 en el que $A_s$ es una extensión de un grupo analítico
abeloide $B_o$ por un toro $T_o$. Utilizamos este toro para definir
$R^t$, de modo que $B$ sea un grupo analítico abeloide sobre $S$.

\medskip
\noindent
\underline{Proposición 13.4.17.} \underline{En las condiciones
precedentes, las filtraciones de Hodge y por pesos sobre}
$\mathcal E=\underline R\otimes_{\mathbb Z}\underline{\mathcal O}_S$,

\begin{equation}\tag{13.4.17.1}
 \mathcal E\supset\underline F,
 \qquad
 \mathcal E\supset\mathcal E^f\supset\mathcal E^t\supset0,
\end{equation}

\noindent
\underline{satisfacen las condiciones siguientes:}

\smallskip
\noindent
a) \underline{La filtración por pesos se obtiene, mediante producto
tensorial sobre $\mathbb Z$ con $\underline{\mathcal O}_S$, de la
filtración análoga sobre $\underline R$}

\begin{equation*}
 \underline R\supset\underline R^f\supset\underline R^t\supset0,
\end{equation*}

\noindent
\underline{donde $\underline R^f$ y $\underline R^t$ son grupos
analíticos localmente constantes sobre} $S$, \underline{con valores en
$\mathbb Z$-módulos libres de tipo finito}, a saber,
$R_o^f=R_s=\pi_1(A_s)$ y $R_o^t=\pi_1(T_o)$.

\smallskip
\noindent
b) \underline{Se tienen las relaciones}

\begin{equation}\tag{*}
 \underline F\cap\mathcal E^t=0,
 \qquad \underline F+\mathcal E^f=\mathcal E\qquad .
\end{equation}

\noindent
\underline{En otras palabras, la filtración sobre $\mathcal E^t$
inducida por la filtración de Hodge}, con las convenciones de grados de
(13.4.6), \underline{está concentrada en grado $-1$}, mientras que la
filtración sobre $\mathcal E/\mathcal E^f$ inducida por la filtración de
Hodge está \underline{concentrada en grado $0$}.

\smallskip
\noindent
c) \underline{Considérese sobre}

\begin{equation*}
 \mathcal E^{\mathrm{ab}}
 \stackrel{\mathrm{def}}{=}\mathcal E^f/\mathcal E^t
 \simeq\underline R^{\mathrm{ab}}\otimes_{\mathbb Z}
 \underline{\mathcal O}_S,
 \qquad
 \underline R^{\mathrm{ab}}=\underline R^f/\underline R^t
 \quad\text{(cf. a))},
\end{equation*}

\noindent
\underline{la filtración inducida por la filtración de Hodge sobre
$\mathcal E$}. Está, por tanto, definida por el submódulo
$\underline F^{\mathrm{ab}}$ de $\mathcal E^{\mathrm{ab}}$ que es la
imagen de

\begin{equation*}
 \underline F^f\stackrel{\mathrm{def}}{=}
 \underline F\cap\mathcal E^f\longrightarrow\mathcal E^{\mathrm{ab}}.
\end{equation*}

\noindent
La flecha es un monomorfismo en virtud de la primera relación $(*)$.
\underline{Entonces el par}

\begin{equation*}
 \bigl(\underline R^{\mathrm{ab}},
 \underline F^{\mathrm{ab}}\subset
 \underline R^{\mathrm{ab}}\otimes_{\mathbb Z}
 \underline{\mathcal O}_S=\mathcal E^{\mathrm{ab}}\bigr)
\end{equation*}

\noindent
\underline{es isomorfo al par}
$(\underline R(B),\underline F(B))$
\underline{asociado a un grupo abeloide sobre $S$}, como se explicó
después de (13.4.4). Por tanto,

\begin{equation*}
 B\simeq\mathcal E^{\mathrm{ab}}/
 (F^{\mathrm{ab}}+R^{\mathrm{ab}}),
\end{equation*}

\noindent
donde los símbolos sin subrayar denotan los grupos analíticos
correspondientes sobre $S$.

\medskip
\noindent
\underline{Observación 13.4.18.} En la terminología de Deligne [4], la
conclusión de la proposición precedente, tras restringirse a un abierto
$U$ de $S$ sobre el que $R$ sea localmente constante, afirma que $R$,
dotado de las dos filtraciones (13.4.17.1) sobre
$\mathcal E=\underline R\otimes_{\mathbb Z}\underline{\mathcal O}_S$,
es una «familia analítica de estructuras de Hodge mixtas sobre $U$».
La situación que acabamos de describir sirvió de modelo para las
conjeturas de Deligne sobre la teoría de Néron en dimensiones
superiores, enunciadas en [9, 9.8--9.13].

Hay, sin embargo, un rasgo importante de la teoría algebraica estudiada
en el \S~7 que falta en la descripción de 13.4.17. Supongamos que $A$ es
un grupo abeloide sobre un abierto $U$ complementario de un subconjunto
analítico $Z$ de $S$ denso en ninguna parte---digamos $Z=\{s\}$ cuando
$\dim S=1$---, de modo que, por (13.1.5), $R|U$ sea un sistema local
sobre $U$. Entonces el sistema local $R/R^t$ sobre $U$ debería ser
\underline{constante}; más precisamente, debería ser dual del sistema
local ${R'}^f$, donde $R'$ es el grupo analítico, étale sobre $S$, que
correspondería (como $R$ corresponde a $A$) a un grupo analítico
adecuado $A'$ sobre $S$ tal que $A'|U$ fuera el grupo abeloide dual de
$A|U$. En la notación del \S~5, el análogo algebraico es la dualidad
entre $T_\ell(A_K)/T_\ell(A_K)^t$ y $T_\ell(A'_K)^f$ de (5.5.9), que
se sigue del teorema de dualidad 5.2.

No sé si tal enunciado es válido bajo las solas hipótesis de 13.4.17,
ni siquiera si la acción de $\pi_1(U,t)$ sobre $R_t$ es unipotente para
$t\in U$. Tampoco parece que los métodos ingenuos utilizados en el
presente párrafo permitan demostrar un resultado semejante, ni aun
bajo la hipótesis adicional de que $A|U$ esté dotado de una
«polarización», es decir, de un haz invertible amplio sobre cada fibra.
Esto es, sin embargo, algo que los «especialistas en grupos
semisimples» parecen saber demostrar, al menos cuando $S$ es no singular
de dimensión $1$; véase el informe [9]. En la subsección siguiente
mostraremos que todo funciona bajo hipótesis de algebraicidad relativa
sobre $A$, y relacionaremos al mismo tiempo el resultado con la
construcción del \S~7.
